Our analysis provides a mean-field perspective on why deep low-rank models can preferentially build high-frequency features through training, beyond what is predicted by purely kernel (NTK/NNGP) limits. We study low-rank (channel-mixed) neural networks in the mean-field regime and show how ReLU gating and structured mixing induce systematic channel specialization and localized spike-like partial functions. We derive channel-wise mean-field dynamics, identify a half-space conditioning mechanism that biases backpropagated signals, and connect this bias to a depth-driven sharpening effect when fitting oscillatory targets. We interpret the resulting training dynamics through an optimization-landscape lens with phase-like transitions, and we report empirical evidence that an Adam-to-SGD schedule can accelerate these transitions. On the theory side, we establish global well-posedness of the associated mean-field ODEs (existence and uniqueness) under sub-Gaussian control via a Picard iteration, and we show that low-rank factorizations with frozen right dictionaries preserve global convergence guarantees for three-layer mean-field training when the induced right span is dense (otherwise yielding convergence to the best approximation in the restricted span). Finally, we report a robust spike-learning phenomenon in random-feature low-rank networks: across experiments, depth \(L\) consistently yields approximately \(2L\) emergent localized spikes that appear sequentially during training.